Below is a list of rewrite rules for key abstract data types, syntactic sugar, and some builtin functions. Phases are intended to be executed in order; the post-condition of one phase serves as the pre-condition for the next. Standard functions (like map, flatten, length, etc.) are written as a library in Simile rather than baked-in optimizations.


\subsection{Syntactic Sugar for Bags}
\begin{align}
  \tag{Bag Image}
  R[B]
  &\leadsto
  R^{-1} \circ B
  \\
  \tag{Bag Predicates - Union}
  S \cup T &\leadsto \Bag{x \in dom(S) \cup dom(T)}{n = max(S[x] \cup T[x])}{x \mapsto n}
  \\
  \tag{Bag Predicates - Intersection}
  S \cap T &\leadsto \Bag{x \in dom(S)}{x \in dom(T) \land n = min(S[x] \cup T[x]) \land n > 0}{x \mapsto n}
  \\
  \tag{Bag Predicates - Sum}
  S + T &\leadsto \Bag{x \in dom(S) \cup dom(T)}{n = sum(S[x] \cup T[x])}{x \mapsto n}
  \\
  \tag{Bag Predicates - Difference}
  % S - T &\leadsto \Bag{x}{x \in dom(S) \land pop(S[x]) - popDefault(T[x], 0) \land n > 0}{x \mapsto n}
  S - T &\leadsto \Bag{x \in dom(S)}{n = S(x) - (T \cup \Bag{x \mapsto 0})(x) \land n > 0}{x \mapsto n}
\end{align}

\subsection{Syntactic Sugar for Sequences}
\begin{align}
    \tag{Concatenation}
    L \concat M
    &\leadsto
    L \cup \Set{x \mapsto y \in M }{x + max(dom(L)) \mapsto y}
\end{align}

\subsection{Syntactic Sugar for Relations}
\noindent\begin{minipage}{\linewidth}
\begin{align}
  \tag{Functional Override}
  f(x) := E
  &\leadsto
  f := f \oplus \Set{x \mapsto E}
  \\
  \tag{Override}
  R \oplus Q &\leadsto Q \cup (dom(Q) \domsub R)
  \\
  \tag{Domain Restriction}
  S \triangleleft R &\leadsto \Set{x \mapsto y \in R}{x \in S}{x \mapsto y}
  \\
  \tag{Domain Subtraction}
  S \domsub R &\leadsto \Set{x \mapsto y \in R}{x \notin S}{x \mapsto y}
  \\
  \tag{Range Restriction}
  R \triangleright S &\leadsto \Set{x \mapsto y \in R}{y \in S}{x \mapsto y}
  \\
  \tag{Range Subtraction}
  R \ransub S &\leadsto \Set{x \mapsto y \in R}{y \notin S}{x \mapsto y}
\end{align}
\end{minipage}

\subsection{Comprehension Construction}

\paragraph{Intuition} All set-like variables and literals are decomposed into set comprehensions.

\paragraph{Post-condition}
\begin{itemize}
  \item All terms with a set-like type (relations, bags, sets, sequences, etc.) must be in comprehension form.
\end{itemize}


\noindent\begin{minipage}{\linewidth}
\begin{align}
  \tag{Functional Image \footnote{May fail if R is not total. $choice$ is nondeterministic.}}
  R(x) &\leadsto choice(R[x])
  \\ % dont use pop - random choice but pop(S) == pop(S). Python op is completely nondeterministic so it may not return the same one R(x) = small epsilon y . y in R[{x}]
  % if R[{x}] is empty, you return "an arbitrary element of range of R", not None
  \tag{Image}
  R[S] &\leadsto \Set{x \mapsto y \in R}{x \in S}{y}
  \\
  \tag{Product}
  x \mapsto y \in S \times T &\leadsto x \in S \land y \in T
  \\
  \tag{Inverse}
  x \mapsto y \in R^{-1} &\leadsto y \mapsto x \in R
  \\
  \tag{Composition}
  x \mapsto y \in (Q \circ R) &\leadsto x \mapsto z \in Q \land z' \mapsto y \in R \land z=z'
  % \tag{Exists Elimination}
  % \exists z \cdot x \mapsto z \in Q \land z \mapsto y \in R &\leadsto x \mapsto z \in Q \land z' \mapsto y \in R \land z=z'
  % \\
  \end{align}
\end{minipage}
\noindent\begin{minipage}{\linewidth}
\begin{align}
  \tag{Predicate Operations - Union}
  S \cup T
  &\leadsto
  \Set{(x \in S \mid x) ` (x \in T \cdot x \notin S \mid x)}
  \\
  \tag{Predicate Operations - Intersection}
  S \cap T
  &\leadsto
  \Set{x \in S}{x \in T}
  \\
  \tag{Predicate Operations - Difference}
  S \setminus T
  &\leadsto
  \Set{x \in S}{x \notin T}
  % \\
  % \tag{Singleton Membership \footnote{Currently unused. We need to be careful to handle the case where $x$ is a free variable.}}
  % x \in \Set{e}
  % &\leadsto
  % x = e
  \\
  \tag{Membership Collapse \footnote{Rule only matches inside the predicate of a quantifier. Explicitly enumerating all matches for all quantuantification types and predicate cases (ANDs, ORs, etc.) would require too much boilerplate. $x$ must be bound by the encasing quantifier.\\
  The $\oplus$ operator represents any quantifier that returns a set-like type (ex. generalized union/intersection, set comprehension, relation comprehension).}}
  x \in \oplus(E \mid P)
  &\leadsto
  P \land x = E
  \\
  \tag{Set Membership}
  e \in S
  &\leadsto
  \exists{x \in S \mid e = x}
\end{align}
\end{minipage}

\subsection{Disjunctive Normal Form}

\paragraph{Intuition} All quantifier predicates are expanded to DNF (i.e. $\land$-operations nested within top-level $\lor$-operations).

\paragraph{Notes} All matches of this phase occur only inside quantifier predicates.

\paragraph{Post-condition}
\begin{itemize}
  \item All terms with a set-like type (relations, bags, sets, sequences, etc.) must be in comprehension form.
  \item Quantifier predicates are in disjunctive normal form - top level or-clauses with inner and-clauses.
  \item If no $\lor$ operators exist within a quantifier predicate, the predicate must only contain $\land$ operators.
\end{itemize}

% https://en.wikipedia.org/wiki/Disjunctive_normal_form
\noindent\begin{minipage}{\linewidth} % Need minipage for footnote
\begin{align}
  \tag{Flatten Nested $\land$}
  x_1 \land ... \land (x_i \land x_{i+1}) \land ...
  &\leadsto
  x_1 \land ... \land x_i \land x_{i+1} \land ...
  \\
  \tag{Flatten Nested $\lor$}
  x_1 \lor ... \lor (x_i \lor x_{i+1}) \lor ...
  &\leadsto
  x_1 \lor ... \lor x_i \lor x_{i+1} \lor ...
  \\
  \tag{Double Negation}
  \lnot \lnot x
  &\leadsto
  x
  \\
  \tag{Distribute De Morgan - Or}
  \lnot (x \lor y)
  &\leadsto
  \lnot x \land \lnot y
  \\
  \tag{Distribute De Morgan - And}
  \lnot (x \land y)
  &\leadsto
  \lnot x \lor \lnot y
  \\
  \tag{Distribute $\land$ over $\lor$}
  x \land (y \lor z)
  &\leadsto
  (x \land y) \lor (x \land z)
\end{align}
\end{minipage}

\subsection{Or-wrapping}
\paragraph{Intuition} Restructuring quantifier predicates with top-level ORs for later rules.

\paragraph{Post-condition}
\begin{itemize}
  \item All terms with a set-like type (relations, bags, sets, sequences, etc.) must be in comprehension form.
  \item Quantifier predicates are in disjunctive normal form - top level or-clauses with inner and-clauses.
\end{itemize}

\noindent\begin{minipage}{\linewidth}
\begin{align}
  \tag{Or-wrapping \footnote{To simplify the matching process later on, we wrap every top-level AND statement (which is guaranteed to be a ListOp by the dataclass field type definition) with an OR.}}
  \Set{E}{\bigwedge P_i}
  &\leadsto
  \Set{E}{\bigvee \bigwedge P_i}
\end{align}
\end{minipage}

\subsection{Generator Selection}
\paragraph{Intuition} All nested ANDs are placed in a tree structure to better suit loop lowering.

\paragraph{Post-condition}
\begin{itemize}
  \item All terms with a set-like type (relations, bags, sets, sequences, etc.) must be in comprehension form.
  \item Each top-level or-clause within quantification predicates must have one selected `generator' predicate (GeneratorSelectionPredicate - GSP) of the form $x \in S$ that loops over its bound dummy variable $x$.
  \item All conjunctive clauses are wrapped in either a GSP or CombGSP structure.
  \item Quantifier predicates are in disjunctive normal form (with GSP replacing AND structures) - top level or-clauses with inner and-clauses.
\end{itemize}

\paragraph{Notes} All matches of this phase occur only inside quantifier predicates. GSP is a tuple of form (generator chain, predicates) that can be flattened to a conjunctive clause. A CombGSP is a triple of form (shared generator, disjunctive children predicates/generators, predicates) that can likewise be flattened into a conjunctive clause. We keep the structure of these tuples distinct to ease the rewriting process.

\paragraph{Brainstorming selection heuristics}
Conditions to consider:
\begin{itemize}
  \item Prefer composition chains in case of nested loops (ex. $x \mapsto y \in R \land y \mapsto z \in Q$, but what about renaming? - $x \mapsto y \in R \land y' \mapsto z \in Q \land y = y'$). Perhaps this could just be a preference for relations over sets if we see a nested quantifier.
  \item Choose most common generator among a list of or-clauses (allows for greater simplification later on)
  \item Choose smallest generator (by set size). This information may not always be accessible (ex. if a set is created by reading values from standard input). Functions may need to choose this on a case-by-case basis (ie. one function call could have two args, with a small set on the left arg and a large set on the right arg. But what if the sizes are reversed later in the code? We've already statically lowered it).
  \item When should constant folding happen? Equality substitution necessary for this?
  \item This may need to work with nesting considerations.
  \item What if we try moving these optimizations to $loop$ structures far later in the pipeline?
\end{itemize}

\noindent\begin{minipage}{\linewidth}
\begin{align}
  \tag{GSP Wrapping \footnote{GSP is a tuple-like structure, where all entries can be flattened into an AND. The LH term must occur inside a quantifier's predicate - one generator per top-level or-clause. Chained generators may be necessary to allow for nested loop operations (ie. non-top-level clauses may require generators). $P_g$ is a list of chained generators, specific set membership clauses (of form $x \in S$) distinguished from the rest of $\bigwedge P_i$. Currently, the selection of $P_g$ is arbitrary (and thus the rewrite system is not confluent), but heurestics may be added later to choose better generators.}}
  \bigwedge P_i
  &\leadsto
  GSP(\bigwedge P_{g_i}, \bigwedge_{P_i \notin P_g} P_i)
  \\
  \tag{Nested Generator Selection \footnote{$free$ is a function that returns true when variables in the predicate are defined (either bound by the current generator or bound prior to the current statement)}}
  GSP(x \land xs, P) \lor GSP(x \land ys, Q)
  &\leadsto
  CombGSP(x, GSP(xs, P) \lor GSP(ys, Q))
\end{align}
\end{minipage}



\subsection{GSP to Loop}

\paragraph{Intuition} Start lowering expressions into imperative-like loops.

\paragraph{Post-condition}
\begin{itemize}
  \item All quantifiers are transformed into $loop$ structures.
  \item $loop$ predicates are of the form $x \in S \land \bigwedge P_i$.
  \item All $loop$ predicates have an assigned generator of the form $x \in S$.
\end{itemize}

\noindent\begin{minipage}{\linewidth}
\begin{align}
  \tag{Quantifier Generation \footnote{$\oplus$ works for any quantifier (but not $\forall$ and $\exists$). The identity and accumulate functions are determined by the realized $\oplus$. For example, if $\oplus = \sum$, the identity is 0 and accumulate is addition.}}
  \oplus E \mid P
  &\leadsto
  \begin{minipage}[]{0.4\textwidth}
  \begin{algorithmic}
  \State $a := identity(\oplus)$
    \Loop\ $P:$ % should we use the \in operator here, or just plaintext in?
        \State $a := accumulate(a, E)$
    \EndLoop
  \end{algorithmic}
  \end{minipage}
\end{align}
\end{minipage}
% \noindent\begin{minipage}{\linewidth}
% \begin{align}
%   \tag{Nested Quantifier}
%   \begin{minipage}[]{0.5\textwidth}
%   \begin{algorithmic}
%     \State $a := accumulator(a, \oplus(E \mid P))$
%   \end{algorithmic}
%   \end{minipage}
%   &\leadsto
%   \begin{minipage}[]{0.4\textwidth}
%   \begin{algorithmic}
%     \Loop\ $P:$
%       \State $a := accumulator(a, E)$
%     \EndLoop
%   \end{algorithmic}
%   \end{minipage}
% \end{align}
% \end{minipage}
\noindent\begin{minipage}{\linewidth}
\begin{align}
  \tag{Top-level Or-Loop}
  \begin{minipage}[]{0.3\textwidth}
  \begin{algorithmic}
    \Loop\ $\bigvee P_i:$
      \State body
    \EndLoop
  \end{algorithmic}
  \end{minipage}
  &\leadsto
  \begin{minipage}[]{0.5\textwidth}
  \begin{algorithmic}
    \Loop\ $P_0$
      \State $body$
    \EndLoop
    \Loop\ $P_1 \land \lnot P_0$
      \State $body$
    \EndLoop
    \Loop\ $P_2 \land \lnot \bigvee_{i < 2} P_i$
      \State $body$
    \EndLoop
  \end{algorithmic}
  \end{minipage}
\end{align}
\end{minipage}
%
\noindent\begin{minipage}{\linewidth}
\begin{align}
  \tag{Chained GSP Loop \footnote{$free$ considers defined variables at this point in time.}}
  \begin{minipage}[]{0.4\textwidth}
  \begin{algorithmic}
    \Loop\ $GSP(g \land gs, P)$
      \State body
    \EndLoop
  \end{algorithmic}
  \end{minipage}
  &\leadsto
  \begin{minipage}[]{0.4\textwidth}
  \begin{algorithmic}
    \Loop\ $SingleGSP(g, free(P))$
      \Loop\ $GSP(gs, bound(P))$
        \State body
      \EndLoop
    \EndLoop
  \end{algorithmic}
  \end{minipage}
\end{align}
\end{minipage}
%
% \noindent\begin{minipage}{\linewidth}
% \begin{align}
%   \tag{Single GSP Loop}
%   \begin{minipage}[]{0.4\textwidth}
%   \begin{algorithmic}
%     \Loop\ $GSP(g, P)$
%       \State body
%     \EndLoop
%   \end{algorithmic}
%   \end{minipage}
%   &\leadsto
%   \begin{minipage}[]{0.4\textwidth}
%   \begin{algorithmic}
%     \Loop\ $SingleGSP(g, P)$
%       \State body
%     \EndLoop
%   \end{algorithmic}
%   \end{minipage}
% \end{align}
% \end{minipage}
%
\noindent\begin{minipage}{\linewidth}
\begin{align}
  \tag{Empty GSP Loop}
  \begin{minipage}[]{0.4\textwidth}
  \begin{algorithmic}
    \Loop\ $GSP([], P)$
      \State body
    \EndLoop
  \end{algorithmic}
  \end{minipage}
  &\leadsto
  \begin{minipage}[]{0.4\textwidth}
  \begin{algorithmic}
    \If{$P$}
      \State body
    \EndIf
  \end{algorithmic}
  \end{minipage}
\end{align}
\end{minipage}
%
\noindent\begin{minipage}{\linewidth}
\begin{align}
  \tag{Combined GSP Loop}
  \begin{minipage}[]{0.4\textwidth}
  \begin{algorithmic}
    \Loop\ $CombGSP(g, gs, P)$
      \State body
    \EndLoop
  \end{algorithmic}
  \end{minipage}
  &\leadsto
  \begin{minipage}[]{0.4\textwidth}
  \begin{algorithmic}
    \Loop\ $SingleGSP(g, P)$
      \Loop\ $gs_0$
        \State body
      \EndLoop
      \Loop\ $gs_1$
        \State body
      \EndLoop
      \State ...
    \EndLoop
  \end{algorithmic}
  \end{minipage}
\end{align}
\end{minipage}

% \noindent\begin{minipage}{\linewidth}
% \begin{align}
%   \tag{Disjunct conditional}
%   \begin{minipage}[]{0.33\textwidth}
%   \begin{algorithmic}
%     \Loop\ $\bigvee GSP_i(P_{g_i}, P_i)$
%       \State body
%     \EndLoop
%   \end{algorithmic}
%   \end{minipage}
%   &\leadsto
%   \begin{minipage}[]{0.5\textwidth}
%   \begin{algorithmic}
%     \Loop\ $GSP_0(P_{g_0}, P_0):$
%       \State body
%     \EndLoop
%     \Loop\ $GSP_1(P_{g_1}, P_1 \land \lnot GSP_0):$
%       \State body
%     \EndLoop
%     \Loop\ $GSP_2(P_{g_2}, P_2 \land \lnot (\bigvee_{i < 2} GSP_i)):$
%       \State body
%     \EndLoop
%     \State ...
%   \end{algorithmic}
%   \end{minipage}
%   \\
%   \tag{Disjunct Conditional 2}
%   \begin{minipage}[]{0.3\textwidth}
%   \begin{algorithmic}
%     \Loop\ $GSP(P_g, \bigvee P_i)$
%       \State body
%     \EndLoop
%   \end{algorithmic}
%   \end{minipage}
%   &\leadsto
%   \begin{minipage}[]{0.8\textwidth}
%   \begin{algorithmic}
%     \Loop\ $GSP(P_g, \bigwedge True):$
%       \If{$\bigvee_{ P_i \neq GSP(...)} P_i$}
%         \State body
%       \EndIf
%       \Loop\ $\bigvee_{P_i: GSP(P_{ig}, P_{is})} GSP(P_{g_i}, P_{is} \land \lnot \bigvee_{P_i \neq GSP(...)} P_i):$
%         \State body
%       \EndLoop
%     \EndLoop
%   \end{algorithmic}
%   \end{minipage}
% \end{align}
% \end{minipage}

\subsection{Relational Subtyping Loop Simplification}

\paragraph{Intuition} Simplify unnecessary iteration structures into direct accesses.

\paragraph{Post-condition}
\begin{itemize}
  \item All quantifiers are transformed into $loop$ structures.
  \item $loop$ predicates are of the form $x \in S \land \bigwedge P_i$.
  \item All $loop$ predicates have an assigned generator of the form $x \in S$.
\end{itemize}
No change from previous, but the number of loop structures should not increase.

\noindent\begin{minipage}{\linewidth}
\begin{align}
  \tag{Total membership elimination \footnote{$R$ is total and $x$ satisfies the type of $R$'s domain. Currently unused since it may eliminate some generators - move higher in the list?}}
  x \in dom(R)
  &\leadsto
  True
  \\
  \tag{Surjective membership elimination \footnote{$R$ is surjective and $x$ satisfies the type of $R$'s range. Currently unused since it may eliminate some generators - move higher in the list? Or }}
  x \in ran(R)
  &\leadsto
  True
\end{align}
\end{minipage}
\noindent\begin{minipage}{\linewidth}
\begin{align}
  \tag{Concrete Domain Image}
  \begin{minipage}[]{0.55\textwidth}
  \begin{algorithmic}
    \Loop\ $SingleGSP(x \mapsto y \in R, x = a \land P)$
      \State body
    \EndLoop
  \end{algorithmic}
  \end{minipage}
  &\leadsto
  \begin{minipage}[]{0.6\textwidth}
  \begin{algorithmic}
    \If{$a \in dom(R)$}
      \Loop\ $SingleGSP(y \in R[a], P[x := a])$
        \State body
      \EndLoop
    \EndIf
  \end{algorithmic}
  \end{minipage}
\end{align}
\end{minipage}
\noindent\begin{minipage}{\linewidth}
\begin{align}
  \tag{Concrete Range Image \footnote{Requires efficient lookups for $ran$, inverse, and image}}
  \begin{minipage}[]{0.55\textwidth}
  \begin{algorithmic}
    \Loop\ $SingleGSP(x \mapsto y \in R, y = a \land P)$
      \State body
    \EndLoop
  \end{algorithmic}
  \end{minipage}
  &\leadsto
  \begin{minipage}[]{0.6\textwidth}
  \begin{algorithmic}
    \If{$a \in ran(R)$}
      \Loop\ $SingleGSP(x \in R^{-1}[a], P[y := a])$
        \State body
      \EndLoop
    \EndIf
  \end{algorithmic}
  \end{minipage}
\end{align}
\end{minipage}
\noindent\begin{minipage}{\linewidth}
\begin{align}
  \tag{Single Element Loop}
  \begin{minipage}[]{0.5\textwidth}
  \begin{algorithmic}
    \Loop\ $SingleGSP(y \in R[x], P)$
      \State body
    \EndLoop
  \end{algorithmic}
  \text{where $free(x)$ and $R$ is many-to-one}
  \end{minipage}
  &\leadsto
  \begin{minipage}[]{0.5\textwidth}
  \begin{algorithmic}
    \If{$P[y:=R(a)]$}
      \State $body[y:=R(a)]$
    \EndIf
  \end{algorithmic}
  \end{minipage}
\end{align}
\end{minipage}
\noindent\begin{minipage}{\linewidth}
\begin{align}
  \tag{Singleton Membership Elimination}
  \begin{minipage}[]{0.5\textwidth}
  \begin{algorithmic}
    \Loop\ $SingleGSP(x \in \Set{y}, P)$
      \State body
    \EndLoop
  \end{algorithmic}
  \end{minipage}
  &\leadsto
  \begin{minipage}[]{0.5\textwidth}
  \begin{algorithmic}
    \If{$P[x:=y]$}
      \State $body[x:=y]$
    \EndIf
  \end{algorithmic}
  \end{minipage}
\end{align}
\end{minipage}

\subsection{Loop Code Generation}

\paragraph{Intuition} Eliminate all intermediate loop structures.

\paragraph{Post-condition}
\begin{itemize}
  \item AST is in imperative code style (for loops, if statements, etc).
  \item All $loop$ and quantification constructs have been eliminated.
  \item Some variables may not be defined.
\end{itemize}

\noindent\begin{minipage}{\linewidth}
\begin{align}
  \tag{Conjunct Conditional \footnote{Function $free$ returns clauses in $P$ that contain only free + defined variables. $bound$ returns the clauses that contain the bound variable $x$ or undefined variables. $P_g$ is the selected generator.}}
  \begin{minipage}[]{0.4\textwidth}
  \begin{algorithmic}
    \Loop\ $SingleGSP(P_g, \bigwedge P_i)$
      \State body
    \EndLoop
  \end{algorithmic}
  \end{minipage}
  &\leadsto
  \begin{minipage}[]{0.5\textwidth}
  \begin{algorithmic}
    \If{$\bigwedge_{free(P_i)} P_i$}
      \For{$P_g$}
        % \State $P_as$ % TODO rewrite these to be more clear that they are assignments
        \If{$\bigwedge_{bound(P_i)} P_i$}
            \State body
          % \If{$Q_0$}
          %   \State body
          % \EndIf
          % \If{$Q_1 \land \lnot Q_0$}
          %   \State body
          % \EndIf
          % \If{$Q_2 \land \lnot (\bigvee_{i < 2} Q_i)$}
          %   \State body
          % \EndIf
          % \State ...
        \EndIf
      \EndFor
    \EndIf
  \end{algorithmic}
  \end{minipage}
\end{align}
\end{minipage}



\subsection{Replace and Simplify}

\paragraph{Intuition} Eliminate all undefined variables (with implicit $\exists$ quantifiers).

\paragraph{Post-condition}
\begin{itemize}
  \item AST is in imperative code style (for loops, if statements, etc).
  \item All variables are defined.
\end{itemize}

\noindent\begin{minipage}{\linewidth}
\begin{align}
  \tag{Equality Elimination \footnote{$x$ is an undefined, unbound variable in the current scope. $E$ is an expression that does not contain $x$. Simplify resulting booleans.}}
  \begin{minipage}[]{0.5\textwidth}
  \begin{algorithmic}
    \If{$Identifier(x) = E \land P$}
      \State body
    \EndIf
  \end{algorithmic}
  \end{minipage}
  &\leadsto
  \begin{minipage}[]{0.4\textwidth}
  \begin{algorithmic}
    \If{$P[x := E]$}
      \State body[x := E]
    \EndIf
  \end{algorithmic}
  \end{minipage}
\end{align}
\end{minipage}
%
\noindent\begin{minipage}{\linewidth}
\begin{align}
  \tag{Simplify And-true}
  P \land True
  &\leadsto
  P
  \\
  \tag{Simplify And-false}
  P \land False
  &\leadsto
  False
  \\
  \tag{Simplify Or-true}
  P \lor True
  &\leadsto
  True
  \\
  \tag{Simplify Or-false}
  P \lor False
  &\leadsto
  P
  \\
  \tag{Flatten Nested Statements}
  Statement(Statement(x))
  &\leadsto
  Statement(x)
\end{align}
\end{minipage}
